\section{Discussion}
\label{sec:discussion}

本节讨论前述实验结果的含义与边界。按照实验设计，\cref{sec:experiments} 已经完成了“观察到什么”的报告；因此，这里重点回答“这些现象说明了什么”以及“在多大程度上可以这样理解”。

\subsection{What the Backbone Comparison Suggests}

\cref{tab:backbone_baseline_comparison} 显示，backbone 深度与验证性能之间并未呈现简单的单调关系。ResNet-34 相比 ResNet-18 没有带来稳定增益，而 ResNet-50 则在主要指标上取得最优结果。这样的结果提示，在当前 $50 \times 50$ 局部病理 patch 任务中，网络加深的收益更可能表现为非平滑的容量阈值效应（capacity threshold effect），而不是线性累积的深度红利。

一种自然的解释是：当模型容量仍处于较低到中等区间时，额外层数未必足以转化为可稳定利用的判别能力；而当表征能力跨过某个阈值后，更深的残差主干才开始更充分地编码局部病理形态。与此同时，这一解释仍然是经验性的，而非统计上已经封闭的结论，因为当前比较建立在统一协议、固定随机种子与单次正式运行之上。若要更稳健地讨论深度收益是否稳定存在，仍需补充多随机种子重复实验、均值与标准差报告，乃至 patient-level 重采样（resampling）分析。

\subsection{Why the Split Protocol Matters}

\cref{tab:split_protocol_comparison} 的核心价值，不在于判断哪一种协议“分数更高”，而在于揭示不同划分方式会如何改变训练表现与验证表现之间的对应关系。实验结果表明，patch-wise random split 在三种 backbone 上都系统性地扩大了 train-eval gap；这意味着模型在训练阶段更容易呈现乐观信号，但这种乐观信号并没有以同等幅度兑现为验证收益。

从数字病理（digital pathology）的数据结构出发，这一现象是可以理解的。IDC patch 数据天然按患者组织，同一患者内部的 patch 往往共享组织背景、染色风格、扫描条件与局部形态统计。因而，当 patch 在患者内部被随机分配到训练集与验证集时，模型更容易利用患者内相关性（intra-patient correlation）所带来的统计便利。结果不是模型一定在验证集上“全面更高”，而是训练阶段的拟合印象更容易被抬高，进而放大训练与泛化之间的脱节。

因此，所谓 patch-wise 的“虚高”，更准确地说，并不是单指某个验证分数的绝对值偏高，而是指它改变了实验者对模型能力的感知方式。若评估协议本身允许患者内相似样本同时出现在训练与验证两侧，那么模型的表面稳定性就可能部分来自边界宽松，而不是来自对未见患者的真正泛化能力。相对而言，patient-wise split 虽然更严格，却更接近真实应用场景下“未见患者”（unseen patients）的评估边界，因此也更适合作为本文的主协议。

\subsection{Limits of the Current Evidence}

这里也需要明确当前证据的边界，避免把经验观察过早上升为过强结论。第一，当前实验主要是单次正式运行结果，尚不足以支持关于统计显著性（statistical significance）或分布稳定性的强断言。第二，当前讨论集中于 ResNet 系列与统一训练配置，因此结论首先适用于“在本任务、本数据与本配置下，协议变化如何影响结果解释”，而不能自动外推到所有结构、所有训练策略或所有数字病理任务。

这恰恰也是本文方法学立场的一部分：在容易受到数据边界影响的任务中，最重要的不是尽早给出宏大解释，而是先确认实验协议是否足够可信。沿着这一路径，下一步最自然的工作包括：在两种协议下补充多随机种子重复实验，系统报告 train-eval gap 的均值、方差与分布；进一步检查 patient-level 分层重采样的稳定性；以及在保持 patient-wise 边界不变的前提下，再引入更复杂的模型或训练策略进行扩展比较。
